\section{Jets on a terminal variety with trivial canonical class}
\label{sec:geometry}

The reduction to a Calabi--Yau pair leaves a positivity problem for a
nef divisor: full nef dimension must imply bigness. We first treat the
case with trivial canonical bundle and terminal singularities. The
argument follows the scalar and two-slot jet constructions of
\cite[Sections~7--8]{LA}. Its geometric exclusions differ from the
nonvanishing argument there: lower-dimensional instances of $(P_j)$
and the polarized effective birationality theorem replace the
exclusions for a canonical nonvanishing counterexample.

\begin{theorem}\label{thm:terminal-big}
Let $n\ge2$, and assume $(P_j)$ for $j<n$. Let $X$ be a projective
$\Q$-factorial terminal complex variety of dimension $n$ with
$K_X\sim0$. If $L$ is a nef Cartier divisor on $X$ with full nef
dimension, then $L$ is big.
\end{theorem}

The proof occupies this section and Section~\ref{sec:frobenius}.
Assuming that $L$ is not big, we obtain two contrasting jet estimates:
sections on $X$ have small normalized vanishing order, whereas a
projective bundle over $X\times X$ admits a large jet system. The
Frobenius comparison will show that these estimates are incompatible.

\subsection{Normalization and the curve obstruction}

Throughout this section suppose that $X,L$ form a counterexample to
Theorem~\ref{thm:terminal-big}. Fix an ample Cartier divisor $A$ on $X$
and a positive rational number $\epsilon$ such that
\begin{equation}\label{j:epsilon}
 (14\epsilon)^n<\frac14,
 \qquad 80\epsilon<\frac34.
\end{equation}
Let $t>0$ tend to zero along a sequence of rational numbers. Since
$L$ is nef and not big, $L^n=0$. We may therefore choose positive
integers $r=r(t)$ satisfying
\begin{equation}\label{eq:scaling}
 P=P_t=r(L+tA),\qquad
 r\longrightarrow\infty,\qquad P^n\longrightarrow\epsilon^n.
\end{equation}
For example, round $\epsilon/((L+tA)^n)^{1/n}$ to the nearest positive
integer. The divisor $P$ is ample and rational Cartier, and $P-rL$ is
nef. After discarding finitely many terms, we have
\begin{equation}\label{j:volume-window}
 \frac{\epsilon^n}{2}\le P^n\le2\epsilon^n.
\end{equation}
Every curve through a very general point of $X$ has positive integral
$L$-degree, hence $P$-degree at least $r$. Thus a covering family of
curves of bounded $P$-degree will give a contradiction.

We specify how the same obstruction applies to subvarieties and
covers. A resolution $W\to X$ has an effective canonical divisor:
pull back a nowhere-zero canonical form on the smooth locus of $X$
and use terminality. Hence $W$, and therefore $X$, is not uniruled
\cite{BDPP}. By Lemma~\ref{j:very-general-images}, if a smooth
projective variety $Y$ maps generically finitely to a
positive-dimensional subvariety containing a sufficiently very general
point of $X$, then $Y$ is not uniruled and the pulled-back Cartier
divisor $L_Y$ has full nef dimension. In particular, $K_Y$ is
pseudo-effective. If $f=\dim Y<n$, the induction hypothesis $(P_f)$
applied to $(Y,0)$ gives
\begin{equation}\label{j:lower-big}
                         K_Y+L_Y\ \text{big}.
\end{equation}
All occurrences of ``very general'' below include the fixed
conditions of Lemma~\ref{j:very-general-images}.

Consequently, for $f<n$ it suffices to find on such a $Y$ a nef
testing divisor $P_Y$ with
\[
 0<P_Y^f\le C_0,\qquad P_Y-rL_Y\ \text{nef},\qquad
 K_YP_Y^{f-1}\le C P_Y^f,
\]
where $C_0,C$ are independent of $t$. Indeed
$L_YP_Y^{f-1}\le P_Y^f/r$, so Lemma~\ref{lemma:bounded-curves}
applies. We will also use that lemma in dimension $n$, after proving
the required bigness separately. For the rest of this section, a
bound is called uniform if it is independent of $t$; it may depend
on parameters fixed before $t$ tends to zero. Since the dimensions
lie in a finite range, Lemma~\ref{lemma:bounded-curves} excludes all
such configurations once $r$ is sufficiently large.

\subsection{The scalar jet estimate}

Put $\rho=(P^n)^{1/n}$. The first estimate bounds the vanishing order
of every section at a very general point; it is stronger than a
bound on a single fixed linear system.

\begin{lemma}\label{lemma:scalar-bound}
In the setting of \eqref{eq:scaling}, for all sufficiently small $t$
and at a very general point $x\in X$, every nonzero section of $kP$,
where $k$ is a positive integer and $kP$ is Cartier, satisfies
\begin{equation}\label{j:scalar-order}
                         \ord_x(s)\le4k\rho.
\end{equation}
\end{lemma}

\begin{proof}
Suppose that violations occur for arbitrarily small $t$. Fix such a
$t$. At a very general point a violation for some $k$ implies that
the corresponding complete jet kernel is nonzero at the general
point. Indeed the kernel ranks are algebraic rank conditions on the
smooth locus, and there are only countably many choices of $k$ and
of the required integral order. Raising the violating section to
powers allows $k$ to be arbitrarily large and divisible.

Choose $n+1$ levels
\[
 \tau_i=\frac32\rho+i\delta\quad(0\le i\le n),
 \qquad \delta=\frac{2\rho}{n}.
\]
They all lie strictly between $\rho$ and $4\rho$. For a sufficiently
large divisible $k$, their full jet kernels are nonzero and
$k\delta\ge4$. The lowest kernel has a positive-dimensional base
component through its marked point. Otherwise that point would be
an isolated base component of multiplicity at least
$h_0=\lceil k\tau_0\rceil$. The zero-dimensional case of
Lemma~\ref{j:base-component} would give
\[
                1\le (k/h_0)^nP^n<1,
\]
a contradiction.

Lemma~\ref{j:moving-kernel} now supplies a moving family of proper
base components $V$, with dominant incidence and dimension
$1\le f<n$. On a resolution $V^*$, write $P_V,L_V$ for the
restrictions followed by pullback. The lemma gives
\[
 0<P_V^f\le(2/\delta)^{n-f}P^n,
 \qquad K_{V^*}P_V^{f-1}\le cP_V^f,
 \qquad 0<c\le\frac{2(n-f)}{\delta}.
\]
Here the canonical term of the ambient variety vanishes because
$K_X\sim0$. The volume window \eqref{j:volume-window} bounds
$\delta$ above and away from zero, so these are uniform estimates.
The incidence point may be chosen over a very general point of $X$.
Thus \eqref{j:lower-big} gives bigness of $K_{V^*}+L_V$, and
$P_V-rL_V$ is nef. Lemma~\ref{lemma:bounded-curves} contradicts
$r\to\infty$.
\end{proof}

\subsection{The projective bundle and its volume}

We now seek a jet system whose order is much larger than the scalar
bound. The additional direction comes from the fibres of a
projective bundle; its two summands keep track of the two copies of
$X$. For a divisor on $X$, subscripts $1,2$ denote pullback from the
two factors of $X\times X$. Set
\begin{equation}\label{eq:bundle}
 Z=\PP\bigl(\OO_X(L)_1\oplus\OO_X(L)_2\bigr)
       \xrightarrow{\pi}X\times X,
 \qquad \xi=c_1(\OO_Z(1)),
 \qquad R=P_1+P_2+q\xi,
\end{equation}
where $q$ is a positive integer to be fixed shortly. We use the
convention
$\pi_*\OO_Z(a)=\operatorname{Sym}^a(\OO_X(L)_1\oplus\OO_X(L)_2)$
for $a\ge0$. The tautological class $\xi$ is nef, and $R$ is ample.
The product and projective-bundle descriptions show that $Z$ is klt.
The two natural sections, called the axes, have classes
$\xi-L_1$ and $\xi-L_2$. The canonical bundle formula gives
\begin{equation}\label{j:axes-canonical}
                         K_Z\sim-\text{(sum of the two axes)}.
\end{equation}
We call the complement of the axes the torus open.

The fibre of either projection $Z\to X$ over a smooth point is,
up to a constant one-dimensional factor,
\[
 Z_{\mathrm{sl}}=\PP(\OO_X\oplus\OO_X(L)),
 \qquad R_{\mathrm{sl}}=P+q\xi.
\]
Its axes have classes $\xi$ and $\xi-L$, and its canonical divisor
is again minus their sum. The slice is projective and klt. In
particular it has the singularity hypotheses needed to apply the
base-component estimates, using a small $\Q$-factorialization if
necessary as explained after Lemma~\ref{j:moving-kernel}.

Write $d=2n+1$. There are positive constants depending only on $n$
such that, for each fixed $q$ and all sufficiently small $t$,
\begin{equation}\label{j:bundle-volumes}
 c_nq\epsilon^{2n}\le R^d\le C_nq\epsilon^{2n},
 \qquad R_{\mathrm{sl}}^{n+1}\le C_{\mathrm{sl},n}q\epsilon^n.
\end{equation}
To see this, expand in nef classes and use
\[
 \pi_*(\xi^j)=\sum_{a=0}^{j-1}L_1^aL_2^{j-1-a}
 \quad(j\ge1).
\]
The term containing exactly one $\xi$ is
$(2n+1)\binom{2n}{n}q(P^n)^2$, giving the lower bound.
For the upper bound, every occurrence of $L_i$ can be bounded in
mixed nef intersections by $P_i/r$. After factoring out $q$, the
remaining powers of $q/r$ are bounded by one once $r\ge q$.
There are only finitely many terms depending on $n$. The slice
calculation is the same, using $\pi_*(\xi^j)=L^{j-1}$.
Together with \eqref{j:volume-window}, this proves
\eqref{j:bundle-volumes}.

For a nef rational Cartier divisor $R$ and a smooth point $z$, its
\emph{Seshadri constant} is
\[
 \varepsilon(R;z)=\inf_{C\ni z}\frac{R\cdot C}{\mult_z C},
\]
where $C$ ranges over integral curves through $z$. Thus a lower bound
for this constant controls degree relative to multiplicity at the point.

\begin{proposition}\label{prop:large-seshadri}
For a suitable fixed positive integer $q$, and all sufficiently
small $t$ in \eqref{eq:scaling}, there is a smooth point $z\in Z$
in the torus open, lying over smooth points of both factors, such
that
\begin{equation}\label{j:large-seshadri}
                         \varepsilon(R;z)>2n+2.
\end{equation}
More generally, for any prescribed dense open subsets $U_1,U_2$
of the smooth locus of $X$, the point may be chosen over
$U_1\times U_2$. In particular, these opens may be the
isomorphism loci of fixed resolutions of $X$.
\end{proposition}

We prove the proposition in the remainder of this section. Choose a
positive rational gap $\delta$ such that
\begin{equation}\label{j:slice-gap}
             (2/\delta)^nC_{\mathrm{sl},n}\epsilon^n<1.
\end{equation}
Next choose $q$ so large that $d+1$ increasing levels of gap $\delta$,
with lowest level strictly greater than $2n+2$, lie strictly below
$(c_nq\epsilon^{2n})^{1/d}$. These choices precede the limit
$t\to0$.

Fix the prescribed opens $U_1,U_2$, if any. Suppose that
\eqref{j:large-seshadri} fails throughout their smooth torus
preimage for arbitrarily small $t$. For any such $t$, asymptotic
Riemann--Roch and the elementary
jet count at a smooth point show that all the full kernels at the
chosen levels are nonzero for sufficiently large divisible $k$.
Indeed the highest level has $d$-th power smaller than $R^d$.
Failure of \eqref{j:large-seshadri} supplies, through each very
general point, a curve with degree divided by multiplicity smaller
than the lowest level. Every section in the lowest kernel vanishes
identically on that curve: otherwise its intersection with the
curve is at least its vanishing order times the curve's
multiplicity, contradicting the degree inequality. Thus the lowest
base locus has a positive-dimensional component through the mark.
Take $k\delta\ge4$ and apply Lemma~\ref{j:moving-kernel}.

We obtain a family of proper integral components $V\subset Z$ whose
incidence dominates $Z$. For its general members, of dimension $f$,
the estimates are
\begin{equation}\label{j:total-component-bounds}
 \begin{split}
 0<R^f\cdot V&\le(2/\delta)^{d-f}R^d,\\
 K_{V^*}R^{f-1}&\le(K_Z+cR)|_V R^{f-1},
 \qquad 0<c\le2(d-f)/\delta.
 \end{split}
\end{equation}
Here and below classes on a resolution are understood by pullback.
The positivity follows from ampleness of $R$. General members meet
the torus open, so \eqref{j:axes-canonical} makes the ambient
canonical contribution nonpositive. All constants in these bounds
are uniform now that $q,\epsilon,\delta$ have been fixed.

\subsection{Excluding positive-dimensional slice fibres}

We first constrain the two projections of $V$ to $X$. For either
projection, consider the component $F$ of a general fibre of
$V\to\operatorname{im}(V\to X)$ through a general incidence point.
Choose that incidence point in the opens of
Lemma~\ref{j:fiber-restriction}, over very general points of both
factors of $X$, and in the smooth torus open. We claim that
\begin{equation}\label{j:no-intermediate-fibre}
                         \dim F\notin\{1,\ldots,n\}.
\end{equation}

Suppose instead that $1\le f_{\mathrm{sl}}=\dim F\le n$.
The ambient fibre is $Z_{\mathrm{sl}}$, of dimension $n+1$.
The restriction lemma gives a nonzero restricted subseries with
$F$ an isolated base component and base ideal contained generically
in $\mathcal I_F^h$, where $h=\lceil k\delta/2\rceil$.
Generic smoothness is imposed over the actual image of $V$; this is
why a general fibre through the incidence point was chosen.
The point is also smooth in the ambient slice. Applying
Lemma~\ref{j:base-component} on that slice gives
\begin{equation}\label{j:slice-bounds}
 \begin{split}
 0<R_{\mathrm{sl}}^{f_{\mathrm{sl}}}\cdot F
   &\le(2/\delta)^{n+1-f_{\mathrm{sl}}}
                      R_{\mathrm{sl}}^{n+1},\\
 K_{F^*}R_{\mathrm{sl}}^{f_{\mathrm{sl}}-1}
   &\le(K_{Z_{\mathrm{sl}}}+cR_{\mathrm{sl}})|_F
                      R_{\mathrm{sl}}^{f_{\mathrm{sl}}-1},
 \end{split}
\end{equation}
with $c\le2(n+1-f_{\mathrm{sl}})/\delta$. Cartier multiples are
chosen before restriction, so the constant line from the fixed
factor causes no change in this calculation.

Let $W_0\subset X$ be the image of $F$ in the other factor.
There are three cases: $F$ is generically finite over a proper
subvariety of $X$; $F$ contains the projective-line fibres over its
image; or $F$ is a multisection over all of $X$.

\paragraph{A generically finite map to a proper subvariety.}
If $F\to W_0$ is generically finite and $\dim W_0<n$, use
$P_Y=R_{\mathrm{sl}}|_{F^*}$ and the pullback $L_Y$ from $X$.
The first inequality in \eqref{j:slice-bounds} gives a uniform
positive volume bound. Since $F$ meets neither axis generically,
its ambient canonical term is nonpositive. Hence the second
inequality gives a uniform canonical-intersection ratio.
Moreover $P_Y-rL_Y$ is nef. The image $W_0$ contains the very
general projection of the incidence point, so
\eqref{j:lower-big} applies. Lemma~\ref{lemma:bounded-curves}
excludes this case for large $r$.

\paragraph{A projective-line bundle over a proper subvariety.}
If $F\to W_0$ is not generically finite, its general fibre has
dimension one. Since the ambient projection is a projective-line
bundle and $F$ is integral, $F$ is the full bundle over $W_0$.
Write $w=\dim W_0=f_{\mathrm{sl}}-1$.
If $w=0$, then $F$ is a projective-line fibre and
$R_{\mathrm{sl}}\cdot F=q$. The first inequality in
\eqref{j:slice-bounds}, together with
\eqref{j:bundle-volumes} and \eqref{j:slice-gap}, gives
\[
 q\le(2/\delta)^nR_{\mathrm{sl}}^{n+1}
   \le(2/\delta)^n C_{\mathrm{sl},n}q\epsilon^n<q,
\]
which is impossible.

Suppose $w>0$. Resolve $W_0$ to $W_0^*$ and take as a resolution
of $F$ the induced projective bundle
$p_0:F^*\to W_0^*$. Its relative canonical divisor is
$-2\xi+p_0^*L$, exactly the restriction of
$K_{Z_{\mathrm{sl}}}$. Cancelling this term in
\eqref{j:slice-bounds} gives
\begin{equation}\label{j:saturated-canonical}
 p_0^*K_{W_0^*}R_{\mathrm{sl}}^w
                   \le cR_{\mathrm{sl}}^{w+1}\cdot F.
\end{equation}
The divisor $K_{W_0^*}$ is pseudo-effective by the very general
point condition. On expanding the left side and pushing down, its
term with exactly one $\xi$ is
$wqK_{W_0^*}P^{w-1}$; every other nonzero term is nonnegative,
since it pairs the pseudo-effective canonical divisor with nef
classes. On the right, the same expansion used for
\eqref{j:bundle-volumes} gives
\begin{equation}\label{j:saturated-volume}
 (w+1)qP^w\cdot W_0
 \le R_{\mathrm{sl}}^{w+1}\cdot F
 \le C'_nqP^w\cdot W_0,
\end{equation}
using $r\ge q$. The slice degree bound and the first inequality in
\eqref{j:saturated-volume} bound $0<P^w\cdot W_0$ uniformly.
Equations \eqref{j:saturated-canonical}--\eqref{j:saturated-volume}
also give
\[
 K_{W_0^*}P^{w-1}\le\frac{cC'_n}{w}P^w\cdot W_0.
\]
Apply \eqref{j:lower-big} and Lemma~\ref{lemma:bounded-curves}
on $W_0^*$ with the pullbacks of $P,L$. This excludes the second
case.

\paragraph{A multisection over $X$.}
The only remaining possibility is that $f_{\mathrm{sl}}=n$ and
$F\to X$ is generically finite of degree $e>0$. Intersect $F$ with
each of the two Cartier axes and push the resulting cycles down
to $X$. This gives effective integral Weil divisors $D_1,D_2$.
The axes do not contain $F$, and their classes differ by $L$;
projection of the associated rational section, or equivalently its
norm, gives
\begin{equation}\label{j:signed-L}
                         D_1-D_2\sim\pm eL.
\end{equation}
For either axis $S$, the classes $R_{\mathrm{sl}}-qS$ and
$R_{\mathrm{sl}}-P$ are nef. Intersecting with the effective cycle
$S\cdot F$ and then using nefness on $F$ yields
\begin{equation}\label{j:axis-degrees}
 qD_iP^{n-1}\le qS\cdot F\cdot R_{\mathrm{sl}}^{n-1}
                  \le R_{\mathrm{sl}}^n\cdot F.
\end{equation}
Thus both axis-image degrees are uniformly bounded.

Set $G=D_1+D_2$. It is $\Q$-Cartier because $X$ is
$\Q$-factorial. Choose $b>0$ rational sufficiently small that
$(X,bG)$ is klt. We claim
\begin{equation}\label{j:positive-boundary-kappa}
                         \kappa(K_X+bG)>0.
\end{equation}
The Iitaka dimension is at least zero, since
$K_X+bG\sim_\Q bG$. If it were zero, Theorems~\ref{input:models}
and \ref{input:abundance} would give a finite ordinary MMP to a
good model $Y$ with transformed adjoint rationally linearly
trivial. The transform of $K_X$ is linearly trivial as well.
Consequently the effective transform of $G$ is zero: all its
components have been contracted. On a common resolution
$a:U\to X$, $b_Y:U\to Y$, relation \eqref{j:signed-L} then
expresses $a^*L$ up to rational linear equivalence as a signed
$b_Y$-exceptional divisor $E$. Divisors exceptional over $X$ are
also exceptional over $Y$, since the MMP extracts none.
Nefness of $a^*L$ and the negativity lemma imply $E\le0$.
Its class is also pseudo-effective. A nonzero anti-effective
divisor has negative intersection with a sufficiently high power
of an ample divisor, so $E=0$. This contradicts full nef dimension
of $L$, and proves \eqref{j:positive-boundary-kappa}.

The already proved positive-Iitaka-dimension case
(Proposition~\ref{prop:full-positive}) now gives
\[
                         K_X+bG+L\ \text{big}.
\]
Its full-positivity hypothesis holds because $L$ has full nef
dimension and $K_X+bG$ is pseudo-effective.
Take a log resolution $a:Y\to X$ of $G$ and let $D_0$ be the
reduced simple normal crossings divisor consisting of the strict
support of $G$ and every exceptional divisor. Decrease $b$ if
necessary; Proposition~\ref{prop:full-positive} still applies.
Terminality and the support of $a^*G$ then give
\[
 K_Y+D_0+a^*L\ \ge\ a^*(K_X+bG+L)
\]
up to rational linear equivalence and an effective difference.
Indeed the finitely many exceptional discrepancies are positive,
and decreasing $b$ makes each coefficient of $b\,a^*G$ no larger
than the corresponding coefficient of $K_Y-a^*K_X+D_0$.
Hence the left side is big.

Test on $Y$ with $a^*P$. Exceptional divisors contribute zero to
intersection with $(a^*P)^{n-1}$, and $K_X\sim0$. The degree of
the reduced strict support is no larger than $GP^{n-1}$, bounded
by \eqref{j:axis-degrees}; also
$LP^{n-1}\le P^n/r$. Since \eqref{j:volume-window} bounds $P^n$
away from zero, we obtain a uniform constant $C$ such that
\[
 (K_Y+D_0+a^*L)(a^*P)^{n-1}\le C(a^*P)^n.
\]
Lemma~\ref{lemma:bounded-curves}, now with its reduced boundary
$D_0$, gives the final contradiction. This proves
\eqref{j:no-intermediate-fibre}.

Thus a general moving component $V$ has fibre dimension either
zero or $n+1$ over its image in either factor. Fibre dimension
$n+1$ would make $V$ the full inverse image of its image in that
factor. Since $V$ is proper, the image has dimension at most
$n-1$, so the other projection would have fibre dimension between
$1$ and $n$, already excluded. Both projections of $V$ are
therefore generically finite onto their images, and $\dim V\le n$.
If $\dim V<n$, the bounds \eqref{j:total-component-bounds}, the
nef class $R-rL_i$, and \eqref{j:lower-big} again contradict
Lemma~\ref{lemma:bounded-curves}. Only $n$-dimensional
correspondences, dominant over both factors, remain.

\subsection{Correspondences and their branch divisors}

For these remaining $V$, \eqref{j:total-component-bounds} gives
uniform bounds
\begin{equation}\label{j:correspondence-bounds}
 R^n\cdot V\le C,
 \qquad K_{V^*}R^{n-1}\le cR^n\cdot V\le C,
\end{equation}
after increasing $C$. Each projection has bounded degree:
$R-P_i$ is nef, and $P^n$ is bounded below.

For each fixed $t$, resolve the general members in their algebraic
family and shrink the parameter space so that the resolutions form
a smooth projective family with integral fibres and morphisms to
$X$ in both slots. The dominant incidence and all numerical bounds
persist. Finite extensions of the parameter space are allowed.
Over $X_{\mathrm{sm}}$, name the components of the relative
Jacobian divisors as in Lemma~\ref{j:finite-covers}. Consider a
component whose image on a general member is a divisor $J\subset X$.

Such divisors $J$ cannot sweep $X$. If they did, choose a general
one through a very general point. Let $E_J$ be a ramification
divisor over it on the corresponding $V^*$. Pullback of a
canonical form from $X$ gives an effective canonical divisor on
$V^*$. Regularity follows by factoring through a resolution of
$X$ and using its canonical singularities; the coefficient of
$E_J$ is positive because its image meets $X_{\mathrm{sm}}$
generically. The difference $R-P_i$ is nef, so
\eqref{j:correspondence-bounds} and projection formula bound
\[
                    0<d_J:=P^{n-1}\cdot J\le C.
\]
On a resolution $J^*$, Lemma~\ref{j:subadjunction}, applied to
$(X,J)$ at the generic point of $J$, gives
\begin{equation}\label{j:branch-canonical}
 \begin{split}
 K_{J^*}P^{n-2}
 &\le(K_X+J)\cdot J\cdot P^{n-2}\\
 &=J^2P^{n-2}
 \le\frac{(JP^{n-1})^2}{P^n}
 \le C' d_J.
 \end{split}
\end{equation}
The pair is lc at that generic point because $X$ is normal and
hence smooth there. The middle inequality is the Hodge index
inequality, applied on a resolution and obtained by ample
approximation of the pulled-back testing class. It does not require
$J$ to be nef. The final constant is uniform by the volume lower
bound and the bound for $d_J$.

Since $J$ sweeps $X$, its resolution satisfies the very general
point conditions used in \eqref{j:lower-big}. With testing class
$P|_{J^*}$, whose top self-intersection is $d_J$, Equation
\eqref{j:branch-canonical} gives precisely the hypotheses of
Lemma~\ref{lemma:bounded-curves} in dimension $n-1$. This
contradiction excludes sweeping branch divisors.

Lemma~\ref{j:finite-covers} now gives, for each slot and this fixed
$t$, only finitely many possible finite normal Stein covers of
$X$, up to isomorphism over $X$. The point of the preceding step
is that degree bounds alone would not give this conclusion:
non-sweeping branch images provide a fixed smooth branch complement
for the family. Neither that complement nor the resulting finite
lists need be uniform in $t$.

\subsection{Fixed covers and an abelian model}

Each $V^*$ maps birationally to both of its normal Stein covers,
and hence determines a graph of a birational map between a pair
of covers in the finite lists. We explain why some such fixed pair
has a family of graphs dominating its product. Restrict the finite
set of pairs to those realized by at least one original member
$V^*$ satisfying \eqref{j:correspondence-bounds}. The original
incidence dominates $Z$, so the projected incidences for these
pairs cover a dense open of $X\times X$. For a fixed pair, graphs of
birational maps lie in countably many Hilbert schemes. Stratify
these schemes so that the fibres are integral and flat and both
projections are birational; the latter conditions may be tested
by dominance and degree one after shrinking. Because $\C$ is
uncountable, a dense open of $X\times X$ cannot be covered by
countably many proper closed subsets. Thus, for one fixed pair
$T_1,T_2$, some integral graph family has dominant evaluation to
$X\times X$. Since $T_1\times T_2$ is finite over $X\times X$
and has the same dimension, this evaluation is dominant onto
$T_1\times T_2$ as well.

Both covers are non-uniruled because they are finite over $X$.
The birational-group theorem of Hanamura
\cite[Theorems~2.1--2.2]{Hanamura}, in the form recalled in
\cite[Proposition~3.7 and Theorems~3.8--3.9]{Blanc}, implies that
they are birational to an abelian variety. Here are the details
of the application. Identify the two covers birationally by one
map in the family. On a suitable smooth projective birational
model of the resulting non-uniruled variety, the reduced
birational group is a group scheme locally of finite type, with
identity component an abelian variety acting regularly. Conjugate
the graph family to this model and shrink again to obtain a flat
family of integral graphs. By the flat-graph representability
statement, its reduced integral parameter space maps to the
reduced birational group. Its connected image lies in one
translate of the identity component. Product-dominant evaluation
therefore gives a dense orbit for that abelian identity component.
The orbit is a quotient of an abelian variety by a stabilizer,
hence is itself an abelian variety. The fixed cover is birational
to it. This is the birational-group part of the fixed-cover
argument in \cite[Proposition~6.11]{LA}; no canonical
nonvanishing conclusion from that proposition is being used.

This structural conclusion depends only on the two covers. We may
therefore return to an original member $V^*$ belonging to this pair,
which retains \eqref{j:correspondence-bounds}, and resolve
further to obtain a birational morphism
\[
                         p_0:V^*\longrightarrow A_0
\]
to an abelian variety. Higher resolutions do not change the
canonical intersection with the pulled-back $(n-1)$-st power of
$R$. Let $L_0$ be the pullback of $L$ from either slot. It is nef,
Cartier, and of full nef dimension. We claim
\begin{equation}\label{j:abelian-big}
                         K_{V^*}+L_0\ \text{big}.
\end{equation}

The pushforward $N_0=(p_0)_*L_0$ is pseudo-effective on $A_0$,
and is therefore nef. Indeed translations make the
pseudo-effective and nef cones of an abelian variety coincide.
If $N_0$ were not big, the standard description of nef line
bundles on abelian varieties would give a positive-dimensional
abelian subvariety on whose translates $N_0$ is numerically
trivial. One can see this from the semipositive Hermitian form of
the numerical class: its kernel is rational for the lattice,
since the alternating form of the integral Cartier class is
integral, and therefore defines an abelian subvariety.
A general translate meets the image of the exceptional locus of
$p_0$ in codimension at least two, if at all. Complete-intersection
curves in that translate can thus be chosen through very general
points and avoiding this image. Their lifts have $L_0$-degree
zero, contradicting full nef dimension. Hence $N_0$ is big.

By the negativity lemma,
\[
                         p_0^*N_0-L_0=E_0\ge0
\]
is exceptional. The canonical divisor of $V^*$ has an effective
representative $E_K$ with positive coefficient along every
$p_0$-exceptional divisor, because the target is smooth with
trivial canonical bundle. Choose a rational $0<u<1$ sufficiently
small that $E_K-uE_0\ge0$. Then
\[
 K_{V^*}+uL_0\sim_\Q u p_0^*N_0+(E_K-uE_0)
\]
is big. Adding the nef divisor $(1-u)L_0$ proves
\eqref{j:abelian-big}.

We may now apply Lemma~\ref{lemma:bounded-curves} in dimension
$n$ on $V^*$, with testing class $R$ and empty boundary.
The bounds \eqref{j:correspondence-bounds} give both the top
intersection bound and the canonical ratio, while $R-rL_0$ is
nef and \eqref{j:abelian-big} supplies bigness. This final
contradiction excludes the correspondence case.

\paragraph{Conclusion of the proof of
Proposition~\ref{prop:large-seshadri}.}
The alternatives for every moving base component have been
exhausted: first its slice fibres, then components of dimension
less than $n$, sweeping branch divisors, and finally the fixed
covers. In each case the contradiction comes from
Lemma~\ref{lemma:bounded-curves} with constants uniform in $t$.
The finite cover lists were used only for one fixed $t$, and their
sizes do not enter the degree bounds. Thus failures of
\eqref{j:large-seshadri} cannot occur along a sequence $t\to0$.
This proves the proposition, including the additional open
conditions on the chosen point.\hfill$\square$

We have now constructed, under the hypothetical failure of
Theorem~\ref{thm:terminal-big}, the scalar estimate of
Lemma~\ref{lemma:scalar-bound} and the projective-bundle estimate
of Proposition~\ref{prop:large-seshadri}. The next section fixes
one sufficiently small value of $t$ and places these data in the
finite-data Frobenius comparison.
